\section{Metric Definitions}
\label{app:metrics}

This appendix fixes the unit of analysis and denominator for every reported
quantity. A task is valid only when the benchmark environment, persistence step,
and evaluator complete without an infrastructure error. Scenario-level measures
use only complete AppWorld scenario triplets; call-level measures use either all
model proposals or executed actions as stated below. Rejected proposals therefore
remain visible to validity metrics but cannot contribute an environmental side
effect. Rates are computed from counts before rounding and are displayed to two
decimal places. Absolute differences between rates are reported in percentage
points, whereas relative changes are reserved for quantities such as token use.

Completion metrics preserve temporal and causal scope. A proposal records model
intent, a gate decision records runtime authorization, and evaluator success is a
judgment over the persisted final state. One task may make several attempts, be
blocked, gather additional evidence, and later be accepted; attempt sets and task
sets are consequently not interchangeable. Recovery follows the same rule: repair
of a classified failure is a recovery-attempt outcome, while eventual benchmark
success remains a task outcome. These definitions are shared by the tables,
figures, trace exporter, and statistics scripts.

\begin{center}
\begin{minipage}{\linewidth}
\centering
\captionof{table}{Completion and recovery metrics, units, and denominators.}
\label{tab:completion-metric-definitions}
\small
\begin{tabularx}{\linewidth}{@{}lYY@{}}
\toprule
Metric & Definition & Unit / denominator \\
\midrule
TGC & Tasks passing every evaluator test divided by valid tasks & task \\
SGC & Complete scenarios for which all three variants pass divided by complete scenarios & scenario \\
System-accepted completion failure & Evaluator-failed tasks after the system allowed termination divided by tasks whose termination the system allowed & accepted task \\
Model-proposed completion failure & Evaluator-failed tasks with a model completion proposal divided by tasks with such a proposal & proposal task \\
False-block attempt rate & Rejected attempts whose decision-time snapshot already satisfied the contract divided by all rejected attempts & blocked attempt \\
False-block affected tasks & Tasks containing at least one false-block attempt divided by tasks blocked at least once & blocked task \\
Gate rescue rate & Tasks blocked at least once and later successful divided by tasks blocked at least once & blocked task \\
Blocks per rescued task & Blocked attempts among rescued tasks divided by rescued tasks & rescued task \\
Max-turn rate & Tasks reaching the configured turn cap divided by valid tasks & task \\
Verification coverage & Required effects with admissible verification divided by required effects & required effect \\
Recovery-attempt success & Recovery attempts that repair their classified failure divided by triggered recovery attempts & recovery attempt \\
Task success after recovery & Recovered tasks finally passing the evaluator divided by tasks that triggered recovery & task \\
\bottomrule
\end{tabularx}
\end{minipage}
\end{center}

\begin{center}
\begin{minipage}{\linewidth}
\centering
\captionof{table}{Tool reliability and resource metrics, units, and denominators.}
\label{tab:tool-resource-metric-definitions}
\small
\begin{tabularx}{\linewidth}{@{}lYY@{}}
\toprule
Metric & Definition & Unit / denominator \\
\midrule
Invalid-call rate & Invalid proposed calls divided by all model-proposed calls & proposed call \\
Out-of-schema rate & Proposed calls absent from the active schema divided by all model-proposed calls & proposed call \\
Duplicate-call rate & Semantically duplicate executed calls divided by all executed calls & executed call \\
Duplicate-write rate & Duplicate write actions divided by all executed writes & executed write \\
Token statistics & Mean, median, P95, and maximum, separated by agent, predictor, user simulator, judge, and runtime & token / task \\
Latency statistics & Mean, median, P95, and maximum from induction to persistence or safe stop & second / task \\
\bottomrule
\end{tabularx}
\end{minipage}
\end{center}

\paragraph{Completion scopes.}
A \emph{model-proposed completion} is an action proposal and does not imply that
the runtime accepted termination. A \emph{gate-accepted completion} is a runtime
decision based on the evidence available at that step. An
\emph{evaluator-confirmed success} is the benchmark's judgment over the persisted
final state. The three events are recorded independently. Multiple completion
attempts can belong to one task, so attempt-level and task-level denominators are
never interchanged.

For the plain baseline, every \code{complete\_task} proposal is implicitly accepted,
so model-proposed and system-accepted completion-failure rates coincide. For
SafeDesk, blocked proposals remain in the model-level rate but not the
system-accepted denominator. A false-block label is produced only after the run:
the evaluator receives the immutable, redacted database snapshot and serialized
Task Contract captured immediately before the rejected attempt, executes the
standard offline test suite without the agent or runtime, and records whether all
contract-required benchmark predicates were already satisfied. Snapshots that fail
to restore, lack evaluator coverage, or include an unresolved policy confirmation
are excluded and reported separately; they are never silently counted as either
true or false blocks. No evaluator output is available to the online gate.

\paragraph{Call relations.}
Out-of-schema calls are a subset of invalid calls and are not added to invalid
calls. Invalid-call rates use all model-proposed calls, including rejected calls.
Duplicate-call rates use executed calls because a rejected proposal cannot repeat
an environmental effect. An executed call is duplicate only when it has the same
normalized tool name and equivalent normalized arguments as an earlier call, no
intervening relevant state or evidence version changed, and the call is neither
pagination nor an explicit readback verification. Duplicate-write rates use
executed writes. A write is duplicate only when intended effect, target entity,
and critical normalized arguments match and the first write was already observed
or verified. These predicates are computed from Action IR and reducer versions,
not assigned by free-form trace review.

\paragraph{Typed recovery outcomes.}
Recovery success is evaluated against the failure-class exit predicate: schema
recovery requires a validating, executable call; entity recovery requires a unique
valid entity after re-query; environment-mismatch recovery requires readback to
satisfy the expected postcondition; and no-progress recovery requires new
admissible evidence or completion of a previously incomplete subgoal. Repairing
this local predicate does not imply evaluator-confirmed task success.

\paragraph{Resource accounting.}
Agent turns count harness proposal cycles. Total LLM calls additionally include
Predictor, Contract Builder, repair, classifier, replanning, resolver, and
summarizer calls. Agent-proposed tool calls and runtime-generated verification
calls are separate proposal sources; total executed calls counts both after guard
disposition. Total tokens aggregate every model role and are also reported by
role. Wall-clock latency spans contract induction through final persistence or
safe stop.

\paragraph{$\tau^2$-bench.}
Pass$^1$ is the fraction of tasks succeeding in a specified trial. Pass$^4$ is
the fraction succeeding in all \TauTrialsPerTask{} fixed repeated trials. Policy
violation is a tool action that violates a domain rule. Premature tool call is an
action taken before required identity, confirmation, or policy preconditions are
met. Tool-response inconsistency means the final response contradicts observed
tool state. Average turns and agent-only tokens exclude user-simulator and judge
usage.
